\section{Conclusion}\label{sec:conclusion}
We presented a closed-loop quality-triage system for skin-lesion diagnosis on degraded consumer
imagery: per input, it decides when to diagnose directly, when to apply diagnosis-preserving
enhancement, and when to request a retake. Its constructive core is a recipe for enhancing
\emph{safely} --- a feature-level diagnosis-preserving objective whose information loss is bounded by a
mutual-information lower bound (Lemma~\ref{lem:dp}) and confirmed empirically --- around which a
four-channel agent routes by a local, per-band conditional risk guarantee
(Theorem~\ref{thm:agent-compact}). We also chart its limits: enhancement turns net-negative on
extreme contrast loss and on malignant lesions under severe degradation, and global triage does not
yet beat direct diagnosis --- an honest negative we read as motivation for learned routing thresholds.
Mapping where degradation-aware enhancement helps, and where it must yield to re-acquisition, is the
more useful contribution for deploying diagnosis on the quality-varying inputs that curated benchmarks
exclude by design.
